% Pursuing Stacks (cotes 134-2 to 134-8): archive cote 134-2, site page 42 (PDF page 43).
% appendix (letters to Breen). Draft compared: ps1.tex:1203-1212.
% AI-transcribed by Claude Opus 5.5 (claude-opus-5-5, adaptive thinking, effort high), Message Batches API, batch ps1 (msgbatch_01Rv5hCaNTtUt4JHx1yqoHPL), 2026-11-12; not yet checked by eye.
\scanpage{134-2/42}
\mspage{27}

images directes et inverses. Mais du même coup il devient clair que le
contexte ``naturel'' des théorèmes de changement de base en cohomologie étale,
des théorèmes\inserted{du type} de Lefschetz (dits ``faibles'') sur les ``sections hyperplanes'',
tout comme \inserted{de} la notion de profondeur qui y joue un rôle crucial, \hand{doit} \inserted{être} celui
des n-champs. Et que le développement hypothétique de ce contexte\deleted{,} ne ris-%
que pas de se réduire à une jonglerie purement formelle et absolument bor-%
délique avec du ``general nonsense'', mais qu'on se trouvera aussitôt confron-%
té à des tests ``d'utilisabilité'' aussi sérieux que la démonstration des
théorèmes de changement de base et ceux du type de Lefschetz (qui même
dans le contexte commutatif ne sont pas piqués de vers ...). \hand{Idem pour
variantes analytiques complexes etc.}

Je ne sais si ces commentaires te ``passent par dessus la tête'' à ton
tour, ni si elles te donnent l'impression qu'il y aurait peut-être des cho-%
ses intéressantes à tirer au clair. Si cela t'intéresse, je pourrais expli-%
citer sous forme un peu plus systématique quelques ingrédients d'une hypo-%
thétique algèbre homologique non commutative et les liens de celle-ci à
l'algèbre homologique commutative. Plus mystérieux pour moi (et pour cause,
vu mon ignorance en homotopie) seraient les relations entre celle-là et
l'algèbre homotopique, I.e. les structures semi-simpliciales, et je n'ai
que des commentaires assez vagues à faire en ce sens\hand{(*)}.Par ailleurs, je te
rappelle que même l'algèbre homologique \xout{\illegible{1}} commutative n'est pas, il s'en
faut, dans un état satisfaisant, pour autant que je sache, vu qu'on ne sait
\margin{(24)} toujours pas quelle est la ``bonne'' notion de catégorie triangulée\hand{\unsure{(*)}}. Or il me
semble bien clair que ce n'est pas une question purement académique - même
si on a pu se passer de le savoir jusqu'ici (en se bornant comme Monsieur
Jourdain à ``faire de la \unsure{prosé|prose} sans le savoir'' - en travaillant sur des
catégories de complexes, éventuellement filtrés, sans trop se demander quel
-les structures il y a sur ces catégories ...).

Bien cordialement à toi

\hand{A Grothendieck}

\margin{(App.3)} \hand{(*) P.S.} Réflexion faite, j'ai quand même envie de te mettre un peu en ap-%
pétit, en faisant ces ``quelques commentaires assez vagues\hand{''}. Il s'agit du
yoga qu'une \inserted{\deleted{\illegible{1}} (petite)} n-catégorie \inserted{\unsure{en groupoïdes|ou groupoïdes}} (à n-équivalence près \deleted{\unsure{``est}} \inserted{``est} essentiellement la même
chose'' qu'un ensemble semi-simplicial pris à homotopie près et où on négli-%
ge les \hand{$\pi$}$_i$ pour i \hand{$\geq$} n+1 (où, si tu préfères,``où on a tué les groupes d'homo-%
topie en dimension \hand{$\geq$}n+1''). Voici des éléments heuristiques pour ce yoga.
Si K. est un ensemble simplicial (il peut être prudent de le prendre de
Kan) on lui associe une n-catégorie \xout{\unsure{Kx}} C$_n$(K.), dont les 0-objets sont les

\margin{42}
